\documentclass[10pt,a4paper]{amsart}
\usepackage[margin=19mm]{geometry}
\usepackage{amsmath,amssymb,mathtools}
\usepackage[scr=rsfs,cal=euler]{mathalfa}
\usepackage{xfrac}
\usepackage[unicode,hidelinks,bookmarks=false]{hyperref}
\newcommand{\ie}{i.e.\ }
\newcommand{\cf}{cf.\ }
\newcommand{\resp}{respectively\ }
\newcommand{\Subsection}{\subsection}
\providecommand{\nobreakdash}{\nobreak\hbox{-}}
\setlength{\emergencystretch}{2em}
\multlinegap0pt
\usepackage{tikz}
\usetikzlibrary{cd}
\usepackage{amssymb}
\usepackage[all]{xy}
\usepackage{bbm}
\usepackage{float}
\usepackage{caption}
\newtheorem{lemma}{Lemma}
\DeclareMathOperator{\Map}{\sf Map}
\DeclareMathOperator{\Sub}{\sf Sub}
\DeclareMathOperator{\dSub}{\delta{\sf Sub}}
\DeclareMathOperator{\ev}{\mathsf{ev}}
\newcommand{\la}{\leftarrow}
\DeclareMathOperator{\pr}{\mathsf{pr}}
\newcommand{\ra}{\rightarrow}
\newcommand{\w}{\widetilde}
\newcommand{\xra}{\xrightarrow}
\title{A parametrized Pontryagin--Thom theorem}
\author{David Ayala, John Francis}
\date{Source: arXiv:2512.10274v2 (CC BY 4.0)}
\begin{document}
\maketitle

\noindent\emph{Source and licence.} A parametrized Pontryagin--Thom theorem. Version arXiv:2512.10274v2 (CC BY 4.0), distributed by arXiv under the Creative Commons Attribution 4.0 International licence, http://creativecommons.org/licenses/by/4.0/, as recorded in the arXiv licence metadata for this exact version and shown on the abstract page https://arxiv.org/abs/2512.10274v2. Full licence notice: \texttt{licenses/arxiv-2512.10274-CC-BY-4.0.txt}.\par\smallskip

\noindent\textit{Editorial scope.} Counted unit: the article's lemma t75 (source0.tex lines 7113--7118). The article's complete original proof of it is delivered with it (source0.tex lines 7120--7447, one \texttt{proof} environment of 328 lines). No mathematics is written, repaired or re-derived: the statement and the proof are the article's own lines, in the article's own notation, and no retained line is substituted or re-flowed.  Retained uncounted as the source-local context the proof needs: lines 1454--1480; lines 7103--7108.  Nothing was omitted from any retained span, so the package carries no omission record.  No retained line contains a citation, so the wrapper carries no bibliography and no bibliography entry.  The retained lines contain the article's own diagram code, which the wrapper typesets with the standard tikz packages; no image file is delivered and the package declares no asset dependency.  Editorial formatting only, in addition: an AMS \texttt{amsart} preamble that restates the article's own declarations and macros used by the retained lines, namely the environments \texttt{lemma} on separate counters, together with the standard packages those lines need.  The retained environments print in this document as Lemma 1 in order of appearance, whereas the article prints the same items with the numbering of its own full document.  The complete native source of the article as distributed in the arXiv e-print for this exact version is bundled unchanged as \texttt{sources/190810/source0.tex}.
\begin{lemma}\label{lemma.sm.sing}    
Consider a solid diagram
\[\begin{tikzcd}
	{C_0} & T \\
	C
	\arrow["f", from=1-1, to=1-2]
	\arrow[hook', from=1-1, to=2-1]
	\arrow["{\w{f}}"', dashed, from=2-1, to=1-2]
\end{tikzcd}\]
in which $T$ is a manifold, $C$ is a manifold with corners, $C_0 \subset C$ is a closed union of faces, and $f$ is a map whose restriction to each face in $C_0$ is smooth.
\begin{enumerate}
    \item 
    If $T$ is a vector space, then there is a smooth extension $\w{f}$.

    \item 
    If $f$ is proper and $T$ is a vector space whose dimension is greater than that of $C$, then there is a proper smooth extension $\w{f}$.

\item 
If $C_0 \subset C$ is properly acyclic,
then there is a smooth extension $\w{f}$.



\end{enumerate}


\end{lemma}

\begin{equation}\label{e79}
\dSub^k(C)
\longrightarrow
\dSub^k(C_0) \underset{\Sub^k(C_0)} {\times^{\sf h}} \Sub^k(C)
~.
\end{equation}

\begin{lemma}\label{t75}
Let $C$ be a manifold with corners satisfying the property that the closure of each face is a manifold with corners in its own right.  
Let $C_0 \subset C$ be a closed union of faces.
The map~(\ref{e79}) is surjective on path-components.

\end{lemma}

\begin{proof}
Recall that the homotopy pullback of the diagram $\dSub^k(C_0)\ra \Sub^k(C_0) \la \Sub^k(C)$ among topological spaces may be defined as the limit of the diagram among topological spaces
% https://q.uiver.app/#q=WzAsNSxbMCwwLCJYIl0sWzEsMSwiWiJdLFsyLDAsIlpeSSJdLFszLDEsIloiXSxbNCwwLCJZIl0sWzAsMSwiZiIsMl0sWzQsMywiZyJdLFsyLDEsIlxcZXZfMCJdLFsyLDMsIlxcZXZfMSIsMl1d
\[\begin{tikzcd}
	\dSub^k(C_0) && {\Sub^k(C_0)^I} && \Sub^k(C) \\
	& \Sub^k(C_0) && \Sub^k(C_0)
	\arrow[""', from=1-1, to=2-2]
	\arrow["{\ev_0}", from=1-3, to=2-2]
	\arrow["{\ev_1}"', from=1-3, to=2-4]
	\arrow["", from=1-5, to=2-4]
\end{tikzcd}\]
where $\Sub^k(C_0)^I := \Map(I,\Sub^k(C_0))$ is the set of maps from the closed interval $I := [0,1]$ to $\Sub^k(C_0)$, endowed with the compact-open topology, and where $\ev_i$ is evaluation at $i\in \{0,1\}$.
Unpacking this description of the homotopy pullback in the case at hand lends to the following description of a point in the codomain of~(\ref{e79}):
\begin{itemize}
    \item[+] a compact codimension-$k$ submanifold
    \[
    \w{W} \subset C
    ~,
    \]
    
    \item[+] a fiber bundle over $I$ of compact codimension-$k$ submanifolds 
    \[
    W_I \subset I \times C_0
    ~;
    \]

    
\end{itemize}
such that the two subsets of $C_0 = \{0\} \times C_0$ agree:
\begin{itemize}
    \item[-] $\w{W} \cap C_0 = W_I \cap (\{0\} \times C_0)$~.
\end{itemize}
By smooth approximation, we may homotope the continuous map $I \to \Sub^k(C_0)$ classifying $W_I \subset I \times C_0$ to assume $W_I \subset I \times C_0$ is a submanifold.  

Through this description of a point in the codomain of the map~(\ref{e79}), this map is given by 
\[
(\w{W} \subset C)
\longmapsto 
\Bigl( ~(\w{W} \subset C)~,~(I \times (\w{W} \cap C_0))~ \Bigr)
~.
\]
Evidently, this map~(\ref{e79}) is injective, and its image consists of those $(\w{W},W_I)$ for which $W_I = I \times (\w{W} \cap C_0)$.



Now, let $(\w{W} \subset C , W_I \subset I \times C_0)$ be a point in the codomain of~(\ref{e79}).
In what follows, we will construct a compact subspace $\w{W}_I \subset I \times C$ that satisfies the following properties:
\begin{enumerate}

    \item The subspace
    \[
    \w{W}_I
    ~\subset~
    I \times C
    \]
    is a fiber bundle over $I$ of compact codimension-$k$ submanifolds of $I \times C$.

    
    \item The two subsets of $I \times C_0$ agree:
    \[
    W_I
    ~=~
    \w{W}_I \cap (I \times C_0) 
    ~.
    \]

    \item The two subsets of $C = \{0\} \times C$ agree:
    \[
    \w{W}
    ~=~
    \w{W}_I \cap (\{0\} \times C) 
    ~.
    \]
    
\end{enumerate}
Given such a $\w{W}_I \subset I \times C$, then $(\w{W}_I \cap (\{1\} \times C) \subset \{1\} \times C = C) \in \dSub^k(C)$, and $\w{W}_I$ witnesses a path in the codomain of~(\ref{e79}) from its image to the given element $(\w{W} \subset C , W_I \subset I \times C_0)$.
Therefore, the lemma is proved upon constructing such a $\w{W}_I \subset I \times C$.





Consider the normal vector bundle $\nu_{\w{W}}$ of $\w{W} \subset C$.
Through the tubular neighborhood theorem, choose an open neighborhood $\w{W} \subset \w{N} \subset C$ together with a diffeomorphism under $\w{W}$ from the total space $E(\nu_{\w{W}}) \cong \w{N}$.
Similarly, consider the normal vector bundle $\nu_{W_I}$ of $W_I \subset I \times C_0$.
Through the tubular neighborhood theorem, choose an open neighborhood $W_I \subset N_I \subset I \times C_0$ together with a diffeomorphism under $W_I$ and over $I$ from the total space $E(\nu_{W_I}) \cong N_I$.
Denote the compact codimension-$k$ submanifold $W_0 = \w{W} \cap C_0 = W_I \cap (\{0\} \times C_0) \subset C_0$.
Upon shrinking and reparametrizing as needed, we may assume $N_0 := \w{N} \cap C_0 = N_I \cap (\{0\} \times C_0) \subset C_0$ and the composite isomorphism 
\[
E(\nu_{W_0}) 
~\cong~
E(\nu_{\w{W}})_{|W_0}
~\cong~
\w{N} \cap C_0 = N_I \cap (\{0\} \times C_0) 
~\cong~ 
E(\nu_{W_I})_{|W_0} 
~\cong~ 
E(\nu_{W_0})
\]
is the identity.
Consider the compact codimension-$k$ submanifold $I \times \w{W} \subset I \times C$.
These data determine a tubular neighborhood $I \times \w{W} \subset I \times \w{N} \subset I \times C$ together with a diffeomorphism $E(\nu_{I\times \w{W}}) = I \times E(\nu_{\w{W}}) \cong I \times \w{N}$, which is 
under $I \times \w{W}$ and over $I$ from the total space $E(\nu_{W_I}) \cong N_I$.
These data, in turn, determine smooth retractions
\[
\pi_I
\colon 
N_I
\cong
E(\nu_{W_I})
\to
W_I
~,\qquad
\w{\pi}
\colon 
\w{N}
\cong
E(\nu_{\w{W}})
\to
\w{W}
~,\text{ and }\qquad
\pi_0
\colon 
N_0
\cong
E(\nu_{W_0})
\to
W_0
~,
\]
the last of which is a common restriction of each of the first two.







{\bf Containment Case.}
We now prove the lemma in the case that there is containment between subsets of $I \times C_0$,
\[
W_I 
~\subset~ 
I \times N_0 
~,
\]
and the resulting map over $I$,
\[
\pr \colon W_I 
~\hookrightarrow~ 
I \times N_0 
\xra{~I \times \pi_0~}
I \times W_0
~,
\]
is a diffeomorphism.
In this case, the composite map
\begin{equation}\label{x22}
I \times W_0
\xra{~\pr^{-1}~}
W_I
\hookrightarrow
I \times N_0
~\cong~
I \times E(\nu_{W_0})
~=~
E(\nu_{I \times W_0})
\end{equation}
is a smooth section of the normal vector bundle $\nu_{I \times W_0}$ over $I \times W_0 \subset I \times C_0$.
Consider the zero-section
\begin{equation}\label{x23}
\w{W}
\xra{~\rm zero~}
E(\nu_{\w{W}})
~,
\end{equation}
which is a smooth section of the normal vector bundle $\nu_{\w{W}}$ over $\w{W} \subset C$.
Observe that~(\ref{x22}) restricts along $W_0 \hookrightarrow I \times W_0$ as the zero-section.  
So the sections~(\ref{x22}) and~(\ref{x23}), together, define a smooth section 
of $\nu_{I \times \w{W}}$ over the closed union of faces $(I \times W_0) \cup \w{W} \subset I \times \w{W}$:
\begin{equation}\label{x24}
\w{W} \cup W_I
\xra{~(\ref{x22})\cup (\ref{x23})~}
E( \nu_{I \times \w{W}})
~.
\end{equation}
Now, using compactness of $I \times \w{W}$, choose a smooth vector bundle injection $\nu_{I \times \w{W}} \hookrightarrow \epsilon^D_{I \times \w{W}}$ into a high-rank trivial vector bundle over $I\times \w{W}$; choose, additionally, a smooth vector bundle retraction $\epsilon^D_{I \times \w{W}} \to \nu_{I \times \w{W}}$ to this injection.  
Lemma~\ref{lemma.sm.sing} then implies the section~(\ref{x24}) extends as a smooth section:
\begin{equation}\label{x25}
I \times \w{W}
\longrightarrow
E(\nu_{I \times \w{W}})
~.
\end{equation}
We then have a composite smooth embedding:
\begin{equation}\label{x26}
I \times \w{W}
\xra{~(\ref{x25})~}
E(\nu_{I \times \w{W}})
~\cong~
I \times \w{N}
~\subset~
I \times C
~.
\end{equation}
The image of~(\ref{x26}),
\[
\w{W}_I
~:=~
{\sf Image}(\ref{x26})
~\subset~
I \times C
~,
\]
is a compact codimension-$k$ submanifold of $I \times C$.
By construction, $\w{W}_I \cap (I \times C_0) = W_I \subset I \times C_0$ and $\w{W}_I \cap (\{0\} \times C) = \w{W} \subset C$.
Additionally, because $\w{W}_I$ is the image of the section of the normal vector bundle of the (trivial) fiber bundle $I \times \w{W} \subset I \times C$, then $\w{W}_I \subset I \times C$ is a fiber bundle over $I$ of compact codimension-$k$ submanifolds of $C$.
This completes the proof of the lemma in this Containment Case.







{\bf General Case.}
We now prove the lemma in the general case, which will reduce to the Containment Case.
Next, using that $W_I \subset I \times C_0$ is a fiber bundle over $I$ of compact codimension-$k$ submanifolds of $C_0$, choose $0=t_0<t_1<\cdots<t_R=1$ such that, for each $0<k\leq R$, the following conditions are satisfied.
\begin{enumerate}
    \item There containment between subsets of $[t_{k-1},t_k] \times C_0$:
    \[
    W_{[t_{k-1},t_k]}
    ~:=~
    W_I \cap \Bigl( \left[t_{k-1},t_k \right]\Bigr) \times C_0
    ~\subset~
    \left[t_{k-1},t_k\right] \times \Bigl(N_I \cap \Bigl(\left\{t_{k-1} \right\} \times C_0\Bigr) \Bigr)
    ~.
    \]

    \item The resulting composite map over $[t_{k-1},t_k]$,
    \[
    \pr_k
    \colon 
    W_{[t_{k-1},t_k]}
    \hookrightarrow
    [t_{k-1},t_k] \times ( N_I \cap ( \{t_{k-1}\} \times C_0 ) )
    \xra{~[t_{k-1},t_k] \times (\pi_I)_{|}~}
    [t_{k-1},t_k] \times ( W_I \cap ( \{t_{k-1}\} \times C_0 ) )
    \]
    is a diffeomorphism.
\end{enumerate}
We now inductively construct, for each $0\leq k\leq R$, a compact subspace $\w{W}_{[0,t_k]} \subset [0,t_k] \times C$ 
that satisfies the following properties:
\begin{enumerate}

    \item The subspace
    \[
    \w{W}_{[0,t_k]}
    ~\subset~
    [0,t_k] \times C
    \]
    is a fiber bundle over $[0,t_k]$ of compact codimension-$k$ submanifolds of $[0,t_k] \times C$.
    
    \item The two subsets of $[0,t_k] \times C_0$ agree:
    \[
    W_I \cap ([0,t_k] \times C_0)
    ~=~
    \w{W}_{[0,t_k]} \cap ([0,t_k] \times C_0) 
    ~.
    \]

    \item The two subsets of $C = \{0\} \times C$ agree:
    \[
    \w{W}
    ~=~
    \w{W}_{[0,t_k]} \cap (\{0\} \times C) 
    ~.
    \]
    
\end{enumerate}
Note that the case $k=R$ is, then, the sought data.
So completing this inductive construction completes the proof of the lemma.


Necessarily, $\w{W}_{[0,t_0]} = \w{W}$, which supplies the base case of $k=0$ for the inductive construction.
Now assume $k>0$, and such a $\w{W}_{[0,t_{k-1}]} \subset [0,t_{k-1}] \times C$ has been constructed.  
The Containment Case supplies a compact codimension-$k$ submanifold $\w{W}_{[t_{k-1},t_k]} \subset [t_{k-1},t_k] \times C$ that satisfies the following properties:
\begin{enumerate}

    \item The subspace
    \[
    \w{W}_{[t_{k-1},t_k]}
    ~\subset~
    [t_{k-1},t_k] \times C
    \]
    is a fiber bundle over $[t_{k-1},t_k]$ of compact codimension-$k$ submanifolds of $[t_{k-1},t_k] \times C$.

    
    \item The two subsets of $[t_{k-1},t_k] \times C_0$ agree:
    \[
    W_I \cap ([t_{k-1},t_k] \times C_0)
    ~=~
    \w{W}_{[t_{k-1},t_k]} \cap ([t_{k-1},t_k] \times C_0) 
    ~.
    \]

    \item The two subsets of $C = \{t_{k-1}\} \times C$ agree:
    \[
    \w{W}_{[0,t_{k-1}]} \cap ( \{t_{k-1}\} \times C)
    ~=~
    \w{W}_{[t_{k-1},t_k]} \cap (\{t_{k-1}\} \times C) 
    ~.
    \]
    
\end{enumerate}
Denote the union
\[
\w{W}_{[0,t_k]}
~:=~
\w{W}_{[0,t_{k-1}]} \cup \w{W}_{[t_{k-1},t_k]}
~\subset~
[0,t_k] \times C
~.
\]
This subspace $\w{W}_{[0,t_k]} \subset [0,t_k] \times C$ is a compact codimension-$k$ subspace that, by construction, satisfies the three named conditions.


\end{proof}

\end{document}
